% Below Macros for diagram generation generated by Claude on 2026-10-18

% ------------------------------------------------------------
% Unit disks, ellipses, unit spheres, and spheroids/ellipsoids
% (in the style of Lay, Linear Algebra and Its Applications, Sec. 3.3)
% Copy this block into a preamble; it needs only \usepackage{tikz}.
%
%   \EllipseFig[keys]     ellipse x1^2/a^2 + x2^2/b^2 <= 1
%   \UnitDiskFig[keys]    same, with a = b = 1 and intercepts labeled 1
%   \EllipsoidFig[keys]   ellipsoid x1^2/a^2 + x2^2/b^2 + x3^2/c^2 <= 1
%   \SpheroidFig[keys]    ellipsoid with b = a (equatorial radius a, polar c)
%   \UnitSphereFig[keys]  same, with a = b = c = 1 and intercepts labeled 1
%
% Keys (all optional):
%   a, b, c                 semi-axis lengths (numbers, default 1)
%   x1 label, x2 label, x3 label
%                           axis labels (default $x_1$, $x_2$, $x_3$; empty = none)
%   a label, b label, c label
%                           labels at the axis intercepts (default none)
%   name                    region label, e.g. $S$ or $T(S)$ (default none)
%   label above, label below
%                           caption centred over / under the figure's origin,
%                           outside the axes (default none)
%   y extent                make the figure reach at least this far above and
%                           below the origin; in 2D it is also the length of the
%                           x2 half-axis (default: just fit the figure). Give
%                           side-by-side figures the same y extent so their
%                           axes and above/below labels line up.
%   fill                    fill / shading colour (default gray!25)
%   overhang                axis length beyond the intercepts (default 0.6)
%   scale                   tikzpicture scale (default 1.5)
%   azimuth, elevation      3D viewing angles in degrees (default 30, 20)
%   equator=true|false      draw the equator x3 = 0 on 3D figures (default true)
%
% Each figure is its own tikzpicture with its baseline at the origin, so
% figures placed side by side line up along the x1-axis.
% ------------------------------------------------------------
\makeatletter
\newif\iflay@equator
\pgfkeys{
  /lay/.is family, /lay,
  a/.store in=\lay@a,             b/.store in=\lay@b,
  c/.store in=\lay@c,
  x1 label/.store in=\lay@xilabel, x2 label/.store in=\lay@xiilabel,
  x3 label/.store in=\lay@xiiilabel,
  a label/.store in=\lay@alabel,   b label/.store in=\lay@blabel,
  c label/.store in=\lay@clabel,
  name/.store in=\lay@name,
  label above/.store in=\lay@above, label below/.store in=\lay@below,
  y extent/.store in=\lay@yext,
  fill/.store in=\lay@fill,
  overhang/.store in=\lay@over,
  scale/.store in=\lay@scale,
  azimuth/.store in=\lay@az,       elevation/.store in=\lay@el,
  equator/.is if=lay@equator,
  reset/.style={a=1, b=1, c=1,
      x1 label=$x_1$, x2 label=$x_2$, x3 label=$x_3$,
      a label=, b label=, c label=, name=,
      label above=, label below=, y extent=,
      fill=gray!25, overhang=0.6, scale=1.5,
      azimuth=30, elevation=20, equator=true},
}
% \lay@ifnonempty{<macro>}{<code>}: run <code> only if <macro> is nonempty
% (expandable, so it also works in the middle of a TikZ path)
\newcommand{\lay@ifnonempty}[2]{%
  \if\relax\detokenize\expandafter{#1}\relax\else#2\fi}
% Applies y extent and draws label above / label below; must come last in
% the tikzpicture, since it reads the bounding box. The fixed text height
% and depth keep labels on a common baseline (e.g. $D$ next to $E$).
\newcommand{\lay@finish}{%
  \lay@ifnonempty\lay@yext{\path (0,-\lay@yext) (0,\lay@yext);}%
  \coordinate (lay@origin) at (0,0);%
  \lay@ifnonempty\lay@above{%
    \node[anchor=south, yshift=4pt, inner sep=0pt,
      text height=1.6ex, text depth=0.4ex]
    at (current bounding box.north -| lay@origin) {\lay@above};}%
  \lay@ifnonempty\lay@below{%
    \node[anchor=north, yshift=-4pt, inner sep=0pt,
      text height=1.6ex, text depth=0.4ex]
    at (current bounding box.south -| lay@origin) {\lay@below};}}

\newcommand{\EllipseFig}[1][]{%
  \begingroup
  \pgfkeys{/lay, reset, #1}%
  \begin{tikzpicture}[scale=\lay@scale, >=stealth, baseline]
    \pgfmathsetmacro\lay@xmax{\lay@a + \lay@over}
    \pgfmathsetmacro\lay@ymax{\lay@b + \lay@over}
    \lay@ifnonempty\lay@yext{\let\lay@ymax\lay@yext}
    \filldraw[fill=\lay@fill, draw=black]
    (0,0) ellipse [x radius=\lay@a, y radius=\lay@b];
    \draw[->] (-\lay@xmax,0) -- (\lay@xmax,0)
    \lay@ifnonempty\lay@xilabel{node[right] {\lay@xilabel}};
    \draw[->] (0,-\lay@ymax) -- (0,\lay@ymax)
    \lay@ifnonempty\lay@xiilabel{node[above] {\lay@xiilabel}};
    \lay@ifnonempty\lay@alabel{\node[below right] at (\lay@a,0) {\lay@alabel};}
    \lay@ifnonempty\lay@blabel{\node[above left] at (0,\lay@b) {\lay@blabel};}
    \lay@ifnonempty\lay@name{%
      \node[above right] at ({\lay@a*cos(45)},{\lay@b*sin(45)}) {\lay@name};}
    \lay@finish
  \end{tikzpicture}%
  \endgroup}
\newcommand{\UnitDiskFig}[1][]{\EllipseFig[a label=$1$, b label=$1$, #1]}

\newcommand{\EllipsoidFig}[1][]{%
  \begingroup
  \pgfkeys{/lay, reset, #1}%
  \begin{tikzpicture}[scale=\lay@scale, >=stealth, baseline,
      lay label/.style={fill=white, fill opacity=0.7, text opacity=1,
          inner sep=1pt, outer sep=1pt}]
    % (Setup lives inside the tikzpicture so line ends add no stray spaces.)
    % Orthographic view from azimuth/elevation: x1 points toward the
    % viewer (lower left), x2 to the right, x3 up.
    \pgfmathsetmacro\lay@st{sin(\lay@az)}  \pgfmathsetmacro\lay@ct{cos(\lay@az)}
    \pgfmathsetmacro\lay@sp{sin(\lay@el)}  \pgfmathsetmacro\lay@cp{cos(\lay@el)}
    % Silhouette = image of the unit circle under L, where L L^T = P D^2 P^T
    % (P = projection, D = diag(a,b,c)); L is the 2x2 Cholesky factor.
    \pgfmathsetmacro\lay@Sxx{(\lay@a*\lay@st)^2 + (\lay@b*\lay@ct)^2}
    \pgfmathsetmacro\lay@Sxy{\lay@st*\lay@ct*\lay@sp*(\lay@a^2 - \lay@b^2)}
    \pgfmathsetmacro\lay@Syy{(\lay@a*\lay@ct*\lay@sp)^2 + (\lay@b*\lay@st*\lay@sp)^2
      + (\lay@c*\lay@cp)^2}
    \pgfmathsetmacro\lay@Lxx{sqrt(\lay@Sxx)}
    \pgfmathsetmacro\lay@Lyx{\lay@Sxy/\lay@Lxx}
    \pgfmathsetmacro\lay@Lyy{sqrt(\lay@Syy - \lay@Lyx^2)}
    % The front half of the equator is centred on angle alpha.
    \pgfmathsetmacro\lay@alpha{atan2(\lay@st/\lay@b, \lay@ct/\lay@a)}
    % Overhangs in world units that all appear \lay@over long on the page.
    \pgfmathsetmacro\lay@oi{\lay@over/sqrt(\lay@st^2 + (\lay@ct*\lay@sp)^2)}
    \pgfmathsetmacro\lay@oii{\lay@over/sqrt(\lay@ct^2 + (\lay@st*\lay@sp)^2)}
    \pgfmathsetmacro\lay@oiii{\lay@over/\lay@cp}
    \shade[inner color=white, outer color=\lay@fill, variable=\t,
      smooth, samples=73]
    plot[domain=0:360]
      ({\lay@Lxx*cos(\t)}, {\lay@Lyx*cos(\t) + \lay@Lyy*sin(\t)}) -- cycle;
    \draw[variable=\t, smooth, samples=73]
    plot[domain=0:360]
      ({\lay@Lxx*cos(\t)}, {\lay@Lyx*cos(\t) + \lay@Lyy*sin(\t)}) -- cycle;
    \begin{scope}[x={(-\lay@st cm, -\lay@ct*\lay@sp cm)},
        y={(\lay@ct cm, -\lay@st*\lay@sp cm)},
        z={(0 cm, \lay@cp cm)}]
      \iflay@equator
        \draw[variable=\t, smooth, samples=37]
        plot[domain={\lay@alpha-90}:{\lay@alpha+90}]
          ({\lay@a*cos(\t)}, {\lay@b*sin(\t)}, 0);
        \draw[dashed, variable=\t, smooth, samples=37]
        plot[domain={\lay@alpha+90}:{\lay@alpha+270}]
          ({\lay@a*cos(\t)}, {\lay@b*sin(\t)}, 0);
      \fi
      % hidden (inside) parts of the axes are dashed
      \draw[dashed] (0,0,0) -- (\lay@a,0,0);
      \draw[dashed] (0,0,0) -- (0,\lay@b,0);
      \draw[dashed] (0,0,0) -- (0,0,\lay@c);
      \draw[->] (\lay@a,0,0) -- ({\lay@a+\lay@oi},0,0)
      \lay@ifnonempty\lay@xilabel{node[below left] {\lay@xilabel}};
      \draw[->] (0,\lay@b,0) -- (0,{\lay@b+\lay@oii},0)
      \lay@ifnonempty\lay@xiilabel{node[right] {\lay@xiilabel}};
      \draw[->] (0,0,\lay@c) -- (0,0,{\lay@c+\lay@oiii})
      \lay@ifnonempty\lay@xiiilabel{node[above] {\lay@xiiilabel}};
      \lay@ifnonempty\lay@alabel{%
        \node[lay label, below right] at (\lay@a,0,0) {\lay@alabel};}
      \lay@ifnonempty\lay@blabel{%
        \node[lay label, below] at (0,\lay@b,0) {\lay@blabel};}
      \lay@ifnonempty\lay@clabel{%
        \node[lay label, right] at (0,0,\lay@c) {\lay@clabel};}
    \end{scope}
    \lay@ifnonempty\lay@name{%
      \node[above right] at ({\lay@Lxx*cos(45)},
      {\lay@Lyx*cos(45) + \lay@Lyy*sin(45)}) {\lay@name};}
    \lay@finish
  \end{tikzpicture}%
  \endgroup}
\newcommand{\SpheroidFig}[1][]{%
  \begingroup\pgfkeys{/lay, reset, #1}%
  \edef\lay@tmp{\noexpand\EllipsoidFig[#1, b=\lay@a]}%
  \expandafter\endgroup\lay@tmp}
\newcommand{\UnitSphereFig}[1][]{%
  \EllipsoidFig[a label=$1$, b label=$1$, c label=$1$, #1]}
\makeatother